\documentstyle{article}
\begin{document}

\centerline{\bf Computer Algebra and the Method of Modified Equations}

\vskip 0.3cm

\centerline{Mohammad O. Ahmed and Robert M. Corless}
\centerline{Department of Applied Mathematics}
\centerline{University of Western Ontario}
\centerline{London, Ontario N6A 5B7, Canada}
\centerline{Email: rcorless@uwo.ca}

\vskip 0.3cm

The method of modified equations is a venerable technique for the
analysis and understanding of the performance of numerical methods
for the solution of ordinary and partial differential equations.  One
reason it is not too much used in practice, although it gives useful
insight, is the tediousness of the
calculations, which have to be done with every differential equation
and numerical method combination.  We
here discuss a simple Maple program which automates the construction of
modified equations, given a differential equation and a numerical method
we wish to analyze.  The basic idea is that we find another, modified,
differential equation which is a better model of what the numerical
method actually solves.

We give several examples showing that the Maple program can, with little
effort, give us good insight into the behaviour of the numerical method.
In particular, some of the terms in the modified equation have physical
interpretations, which helps the user's intuition.

Our current implementation just uses Taylor series, but work on a
B-series version is under way, and we hope by the time of the conference
to have some preliminary results with which to compare the two
approaches.

\end{document}
